% !TEX program = xelatex
% Slides du TP « sketches pour la comparaison de génomes »
% Thème Metropolis — compiler avec XeLaTeX (police Fira + accents).
% Schémas inspirés du billet de Camille Marchet :
%   https://kamimrcht.github.io/webpage/sketch.html  (autorisation accordée)
\PassOptionsToPackage{table}{xcolor} % active \rowcolor (chargé via beamer)
\documentclass[aspectratio=169,11pt]{beamer}
\usetheme[progressbar=frametitle, numbering=fraction, sectionpage=progressbar, block=fill]{metropolis}

% Police Fira chargée par NOM DE FICHIER via kpathsea (portable : pas de chemin
% absolu, fonctionne sur toute install TeX Live qui fournit le paquet `fira`).
\usepackage{fontspec}
\setsansfont{FiraSans}[
  Extension      = .otf,
  UprightFont    = *-Regular,
  ItalicFont     = *-Italic,
  BoldFont       = *-SemiBold,
  BoldItalicFont = *-SemiBoldItalic,
]
\setmonofont{FiraMono}[
  Extension   = .otf,
  UprightFont = *-Regular,
  BoldFont    = *-Bold,
  Scale       = MatchLowercase,
]

\usepackage[french]{babel}
\usepackage{csquotes}
\usepackage{amsmath,amssymb}
\usepackage{booktabs}
\usepackage{listings}
\usepackage{tikz}
\usetikzlibrary{positioning,arrows.meta,calc,shapes.geometric,decorations.pathreplacing,fit,backgrounds,matrix,babel}

%% ---------------------------------------------------------------- couleurs
\definecolor{cBlue}{HTML}{1F6FB2}
\definecolor{cOrange}{HTML}{EB811B}
\definecolor{cGreen}{HTML}{2E8B57}
\definecolor{cRed}{HTML}{C0392B}
\definecolor{cGray}{HTML}{9AA0A6}
\definecolor{cPurple}{HTML}{6A4C93}

%% ---------------------------------------------------------------- styles tikz
\tikzset{
  arr/.style   = {-{Stealth[length=2.3mm]}, line width=0.5pt},
  nt/.style    = {draw, minimum size=6.5mm, inner sep=0pt, font=\ttfamily\small},
  bit/.style   = {draw, minimum width=6mm, minimum height=6.5mm, inner sep=0pt, font=\ttfamily\footnotesize},
  cell/.style  = {draw, minimum width=7.5mm, minimum height=7.5mm, inner sep=0pt, font=\small},
  bucket/.style= {draw, minimum width=8mm, minimum height=11mm, inner sep=0pt},
}

%% ---------------------------------------------------------------- listings (code)
% Contenu des listings en ASCII pur (évite tout souci d'encodage).
\lstdefinestyle{py}{
  language=Python,
  basicstyle=\ttfamily\scriptsize,
  keywordstyle=\color{cBlue}\bfseries,
  commentstyle=\color{cGreen}\itshape,
  stringstyle=\color{cOrange},
  numberstyle=\tiny\color{cGray},
  showstringspaces=false,
  columns=fullflexible,
  keepspaces=true,
  breaklines=true,
  aboveskip=3pt, belowskip=2pt,
}
\lstset{style=py}

%% ---------------------------------------------------------------- méta
\title{Sketches pour la comparaison de génomes}
\subtitle{k-mers, MinHash, partitions \& HyperMinHash}
\author{TP \textsc{Jc2bimmm}}
\date{}
\institute{Schémas inspirés du billet de C.~Marchet}

\begin{document}

\maketitle

%% ================================================================ MOTIVATION
\begin{frame}{Comparer deux génomes}
  \begin{columns}[T]
    \begin{column}{0.5\textwidth}
      On veut mesurer à quel point deux génomes \alert{se ressemblent}.
      \medskip

      Idée : découper chaque génome en \textbf{k-mers} (mots de $k$ nucléotides),
      puis comparer les \textbf{ensembles} de k-mers.
      \medskip

      \textbf{Indice de Jaccard :}
      \[ J(A,B)=\frac{|A\cap B|}{|A\cup B|}\in[0,1] \]
      \footnotesize Plus $J$ est grand, plus les génomes partagent de k-mers.
    \end{column}
    \begin{column}{0.5\textwidth}
      \centering
      \begin{tikzpicture}[scale=0.92]
        \node[font=\ttfamily\small] (g1) at (0,1.7)  {ACGTTGCA\ldots};
        \node[font=\ttfamily\small] (g2) at (0,-1.7) {ACGATGCA\ldots};
        \node[font=\scriptsize,cGray] at (0,2.15){génome 1};
        \node[font=\scriptsize,cGray] at (0,-2.15){génome 2};
        \begin{scope}[shift={(2.4,0)}]
          \fill[cBlue!25]   (0,0)   ellipse (1.25 and 1.45);
          \begin{scope}\clip (0,0) ellipse (1.25 and 1.45);
            \fill[cGreen!45] (1.35,0) ellipse (1.25 and 1.45);\end{scope}
          \fill[cOrange!25] (1.35,0) ellipse (1.25 and 1.45);
          \begin{scope}\clip (1.35,0) ellipse (1.25 and 1.45);
            \fill[cGreen!45] (0,0) ellipse (1.25 and 1.45);\end{scope}
          \draw (0,0) ellipse (1.25 and 1.45);
          \draw (1.35,0) ellipse (1.25 and 1.45);
          \node at (-0.45,0) {$A$};
          \node at (1.8,0)   {$B$};
          \node[font=\scriptsize] at (0.67,0) {$A\cap B$};
        \end{scope}
        \draw[arr,cGray] (g1.east) to[out=0,in=120]  (2.0,1.1);
        \draw[arr,cGray] (g2.east) to[out=0,in=-120] (2.0,-1.1);
      \end{tikzpicture}
    \end{column}
  \end{columns}
  \vfill
  \footnotesize\textcolor{cRed}{Problème :} un génome a des \emph{millions} de k-mers
  $\Rightarrow$ on veut un \textbf{résumé compact} (un \emph{sketch}).
\end{frame}

%% ================================================================ PARTIE 1
\section{Énumérer les k-mers}

\begin{frame}{Énumérer les k-mers : par sous-chaînes}
  Une séquence de longueur $n$ contient $n-k+1$ k-mers (fenêtre glissante).
  La méthode naïve extrait une sous-chaîne \texttt{seq[i:i+k]}.
  \smallskip

  \centering
  \begin{tikzpicture}[scale=1]
    \foreach \c [count=\i from 0] in {A,C,G,T,T,G,C,A}
      \node[nt] (n\i) at (\i*0.75,1.3) {\c};
    % fenêtre k=3 sur les 3 premiers (A C G)
    \draw[cBlue,line width=1pt,rounded corners] (-0.42,0.9) rectangle (1.86,1.7);
    \draw[arr,cBlue] (0.9,1.95) -- (2.4,1.95)
      node[midway,above,font=\scriptsize,cBlue]{la fenêtre glisse};
    % k-mers extraits
    \node[font=\scriptsize,cGray,anchor=east] at (-0.6,0){k-mers :};
    \foreach \m [count=\j from 0] in {ACG,CGT,GTT,TTG,TGC,GCA}
      \node[font=\ttfamily\footnotesize] at (\j*1.05+0.1,0) {\m};
  \end{tikzpicture}

  \smallskip
  \begin{block}{Complexité (méthode sous-chaînes)}
    \footnotesize
    \begin{itemize}
      \item Extraire un k-mer (copier $k$ lettres) : $O(k)$
      \item Énumérer \emph{tous} les k-mers : $\mathbf{O(n\cdot k)}$ en temps, $O(k)$ mémoire/k-mer
      \item Comparer deux k-mers (chaînes) : $O(k)$ — lettre par lettre
    \end{itemize}
  \end{block}
\end{frame}

\begin{frame}{Encoder un k-mer en entier (2 bits / nucléotide)}
  \begin{columns}[T]
    \begin{column}{0.46\textwidth}
      \footnotesize
      Chaque nucléotide tient sur \textbf{2 bits} :
      \[ \texttt{A}{=}00\quad \texttt{C}{=}01\quad \texttt{G}{=}11\quad \texttt{T}{=}10 \]
      (obtenus par \texttt{(ord(c){>}{>}1)\&0b11})
      \medskip

      Un k-mer $\Rightarrow 2k$ bits, gardé dans un seul entier avec un masque
      \texttt{mask = (1{<}{<}2k)-1}.
    \end{column}
    \begin{column}{0.54\textwidth}
      \centering
      \begin{tikzpicture}
        \foreach \L/\i in {A/0,C/1,G/2,T/3}
          \node[nt] (k\i) at (\i*1.3,0.9) {\L};
        \foreach \B/\i in {00/0,01/1,11/2,10/3}
          \node[bit,fill=cOrange!18] (b\i) at (\i*1.3,0.1) {\B};
        \foreach \i in {0,1,2,3} \draw[arr,cGray] (k\i)--(b\i);
        \draw[decorate,decoration={brace,mirror,amplitude=4pt}]
          (b0.south west)--(b3.south east)
          node[midway,below=5pt,font=\ttfamily\footnotesize]{0b00011110 = 30};
      \end{tikzpicture}
    \end{column}
  \end{columns}
  \medskip
  \begin{alertblock}{La limite des 64 bits}
    \footnotesize
    Un mot machine fait 64 bits $\Rightarrow 2k\le 64 \Rightarrow \mathbf{k\le 32}$.
    Comparer deux entiers $\le 64$ bits coûte \textbf{$O(1)$} (une instruction CPU),
    au lieu de $O(k)$ pour des chaînes.
  \end{alertblock}
\end{frame}

\begin{frame}{Énumérer en fenêtre glissante d'entiers}
  Au lieu de ré-encoder $k$ lettres à chaque position, on \textbf{met à jour} l'entier :
  \[ \texttt{kmer = ((kmer {<}{<} 2) + val) \& mask} \]
  \vspace{-6pt}
  \begin{center}
  \begin{tikzpicture}
    \node[font=\scriptsize,cGray,anchor=east] at (-0.55,0.8){avant :};
    \foreach \L/\i in {A/0,C/1,G/2,T/3} \node[nt] (a\i) at (\i*0.8,0.8) {\L};
    \node[nt,draw=cRed,text=cRed] at (0,0.8) {A};
    \draw[arr,cRed] (a0.north) to[out=120,in=-55] (-1.0,1.5)
      node[left,font=\scriptsize,cRed]{chute (masque)};
    \draw[arr] (1.2,0.45) -- (1.2,-0.15)
      node[midway,right,font=\scriptsize]{$<{<}2$, $+$nt, \&mask};
    \node[font=\scriptsize,cGray,anchor=east] at (-0.55,-0.5){après :};
    \foreach \L/\i in {C/0,G/1,T/2} \node[nt] (c\i) at (\i*0.8,-0.5) {\L};
    \node[nt,fill=cGreen!25] (cnew) at (2.4,-0.5) {A'};
    \node[cGreen,font=\scriptsize,anchor=west] at (3.0,-0.5){$\leftarrow$ nouveau nt};
  \end{tikzpicture}
  \end{center}
  \vspace{-2pt}
  {\footnotesize K-mer \textbf{canonique} : \texttt{min(forward, reverse-complement)}
  $\Rightarrow$ indépendant du brin.}
  \begin{block}{Complexité (fenêtre glissante)}
    \footnotesize
    $O(1)$ par base $\Rightarrow \mathbf{O(n)}$ en temps, $O(1)$ mémoire/k-mer.
    \hfill\textcolor{cGreen}{Gain $\times k$} (naïf $O(n\cdot k)$).
  \end{block}
\end{frame}

%% ================================================================ PARTIE 2
\section{Hachage \& Bottom-$s$ MinHash}

\begin{frame}{Qu'est-ce qu'une fonction de hachage ?}
  \begin{columns}[T]
    \begin{column}{0.54\textwidth}
      \footnotesize
      Une fonction de hachage \texttt{h} prend un entier et en renvoie un autre,
      en \textbf{mélangeant ses bits}. Trois propriétés suffisent :
      \begin{itemize}
        \item \textbf{déterministe} : même entrée $\to$ même sortie ;
        \item \textbf{(quasi-)uniforme} sur $[0,2^b)$ ;
        \item \textbf{casse la proximité} : deux k-mers voisins $\to$ valeurs très éloignées.
      \end{itemize}
      \textcolor{cGray}{On veut de la \emph{diversité}, pas des k-mers qui se ressemblent.}
      \medskip

      \textbf{Une boîte noire :} peu importe \emph{comment} elle calcule.
      Il en existe beaucoup (MurmurHash, xxHash, xorshift\ldots) — on en choisit une.
    \end{column}
    \begin{column}{0.46\textwidth}
      \centering
      \begin{tikzpicture}
        \node[nt] (x1) at (0,0.6) {x};
        \node[nt,minimum width=9mm] (x2) at (0,-0.6) {x+1};
        \node[draw,rounded corners,fill=black,text=white,minimum height=1.7cm,minimum width=1.2cm] (h) at (1.8,0) {$h$};
        \draw[arr] (x1)--(h); \draw[arr] (x2)--(h);
        \draw[-{Stealth}] (0.2,-1.9)--(4.2,-1.9) node[right,font=\scriptsize]{$2^b$};
        \node[font=\scriptsize] at (0.2,-2.2){0};
        \draw[cBlue,line width=1pt]  (1.4,-1.75)--(1.4,-2.05) node[below,font=\scriptsize,cBlue]{$h(x)$};
        \draw[cOrange,line width=1pt](3.4,-1.75)--(3.4,-2.05) node[below,font=\scriptsize,cOrange]{$h(x{+}1)$};
        \draw[arr,cGray] (h.east) to[out=0,in=110] (1.4,-1.7);
        \draw[arr,cGray] (h.east) to[out=0,in=100] (3.4,-1.7);
        \node[font=\scriptsize,cGray] at (1.8,-2.65){une boîte noire};
      \end{tikzpicture}
    \end{column}
  \end{columns}
\end{frame}

\begin{frame}{Bottom-$s$ MinHash : garder les $s$ plus petits hachés}
  \begin{columns}[T]
    \begin{column}{0.5\textwidth}
      \footnotesize
      Sélectionner les \textbf{$s$ plus petits hachés} = tirer un \textbf{échantillon
      aléatoire uniforme} des k-mers, \emph{cohérent} d'un génome à l'autre
      (mêmes k-mers $\to$ mêmes petits hachés).
      \medskip

      \textbf{Principe à coder :} maintenir l'ensemble des $s$ plus petits hachés vus.
      \begin{itemize}
        \item moins de $s$ gardés : on \textbf{ajoute} ;
        \item sinon, si le nouveau $<$ \textbf{plus grand gardé} : il le \textbf{remplace} ;
        \item sinon : on l'\textbf{ignore}.
      \end{itemize}
      \medskip
      Jaccard estimé : $\hat J = \dfrac{|S_A\cap S_B|}{|S_A\cup S_B|}$.
    \end{column}
    \begin{column}{0.5\textwidth}
      \centering
      \begin{tikzpicture}
        \draw[-{Stealth}] (0,0)--(5,0) node[right,font=\scriptsize]{$2^b$};
        \node[font=\scriptsize] at (0,-0.3){0};
        \foreach \x in {0.3,0.6,1.0,1.5,2.2,2.8,3.3,3.7,4.1,4.5}
          \draw[cGray] (\x,0.16)--(\x,-0.16);
        \foreach \x in {0.3,0.6,1.0,1.5}
          \draw[cBlue,line width=1.1pt] (\x,0.22)--(\x,-0.22);
        \draw[decorate,decoration={brace,amplitude=4pt},cBlue]
          (0.25,0.45)--(1.55,0.45)
          node[midway,above,font=\scriptsize,cBlue]{sketch = $s$ plus petits};
        \draw[arr,cRed] (1.5,-0.95)--(1.5,-0.28);
        \node[font=\scriptsize,cRed,align=center] at (1.5,-1.3){plus grand\\gardé};
        \node[font=\scriptsize,cGray,align=center] at (3.6,-0.7){les autres\\sont ignorés};
      \end{tikzpicture}

      \vspace{4pt}
      {\scriptsize\color{cGray} À vous : implémentez \texttt{add\_kmers}.}
    \end{column}
  \end{columns}
\end{frame}

\begin{frame}{Complexité du bottom-$s$ \;(+ bonus : tas)}
  \begin{columns}[T]
    \begin{column}{0.54\textwidth}
      \footnotesize
      \textbf{Version naïve (\texttt{set}) :} recalculer \texttt{max(kept)} à chaque
      remplacement coûte $O(s)$ $\Rightarrow \mathbf{O(n\cdot s)}$ en temps, $O(s)$ mémoire.
      \smallskip

      \textbf{Bonus — tas (\texttt{heapq}) :} on stocke \texttt{-kmer}, le max reste
      à la racine :
      \begin{itemize}
        \item lire le max : $O(1)$
        \item \texttt{heappushpop} (insérer + retirer) : $O(\log s)$
      \end{itemize}
      $\Rightarrow \mathbf{O(n\log s)}$ en temps.
    \end{column}
    \begin{column}{0.46\textwidth}
      \centering
      \begin{tikzpicture}[level distance=7mm,
          level 1/.style={sibling distance=16mm},
          level 2/.style={sibling distance=8mm},
          every node/.style={draw,circle,minimum size=5.5mm,inner sep=0pt,font=\scriptsize}]
        \node[fill=cRed!25] (root) {max}
          child {node {} child {node {}} child {node {}}}
          child {node {} child {node {}} child {node {}}};
      \end{tikzpicture}

      {\scriptsize\color{cGray} tas-max : racine = maximum}

      \smallskip
      {\scriptsize
      \begin{tabular}{lcc}
        \toprule
         & \texttt{set} & tas \\
        \midrule
        max     & $O(s)$ & $O(1)$ \\
        remplacer & $O(s)$ & $O(\log s)$ \\
        \midrule
        \textbf{total} & $O(ns)$ & $O(n\log s)$ \\
        \bottomrule
      \end{tabular}}
    \end{column}
  \end{columns}
\end{frame}

%% ================================================================ PARTIE 3
\section{Partition sketch}

\begin{frame}{Les deux coûts du bottom-$s$}
  \footnotesize
  Les $s$ valeurs gardées dépendent d'un \textbf{ordre global} (les plus petites parmi
  \emph{toutes}) — d'où \textbf{deux} coûts.
  \smallskip

  \textbf{1. Construction (mise à jour).} Un k-mer plus petit que le max \emph{met le
  sketch à jour}, en $O(\log s)$ (réordonner le tas) — potentiellement à \emph{chaque}
  k-mer. \textcolor{cGray}{Sur $n \gg s$ les mises à jour se raréfient ($\Rightarrow$ coût
  \textbf{amorti}), mais pas forcément anodin.}
  \smallskip

  \textbf{2. Comparaison.} Deux sketches n'ont pas les mêmes valeurs aux mêmes positions :
  pour $|S_A\cap S_B|$ il faut \textbf{trier / fusionner} $\Rightarrow \mathbf{O(s\log s)}$.

  \medskip
  \centering
  \begin{tikzpicture}[scale=0.9]
    \node[font=\scriptsize,cGray,anchor=east] at (-0.3,0.6){$S_A$ :};
    \foreach \v/\i in {3/0,9/1,12/2,20/3} \node[cell,fill=cBlue!12] at (\i*1.0,0.6){\v};
    \node[font=\scriptsize,cGray,anchor=east] at (-0.3,-0.6){$S_B$ :};
    \foreach \v/\i in {5/0,9/1,11/2,20/3} \node[cell,fill=cOrange!12] at (\i*1.0,-0.6){\v};
    \node[draw,rounded corners,minimum height=1.5cm,fill=cGray!12] (m) at (4.6,0){fusion};
    \draw[arr] (3.55,0.6) -- (m.west|-0,0.3);
    \draw[arr] (3.55,-0.6) -- (m.west|-0,-0.3);
    \draw[arr] (m.east) -- +(0.9,0) node[right,font=\scriptsize]{$|S_A\cap S_B|$};
  \end{tikzpicture}

  \smallskip
  {\footnotesize\textcolor{cGreen}{$\Rightarrow$ on aimerait des cases \textbf{indépendantes} :
  ni tas à réordonner, ni tri.}}
\end{frame}

\begin{frame}{Partition sketch de taille $s$}
  \footnotesize
  \textbf{Idée :} rendre les $s$ valeurs \emph{indépendantes}. On découpe l'espace
  en $s$ \textbf{partitions} ; chaque k-mer va dans la partition \texttt{p = kmer \% s} ;
  dans chaque partition on ne garde que le \textbf{minimum}.
  \medskip

  \centering
  \begin{tikzpicture}
    % valeurs en entrée
    \foreach \v/\x in {12/0, 7/1, 21/2, 3/3, 15/4, 8/5, 30/6}
      \node[font=\scriptsize] (v\x) at (\x*1.15,1.7) {\v};
    \node[font=\scriptsize,cGray,anchor=east] at (-0.5,1.7){hachés :};
    % 5 partitions
    \foreach \i in {0,...,4}
      \node[bucket] (p\i) at (\i*1.4+0.6,0) {};
    \foreach \i in {0,...,4}
      \node[font=\scriptsize,cGray] at (\i*1.4+0.6,-0.85){\texttt{\i}};
    % minima gardés (12%5=2->? ; on illustre simplement un min par seau)
    \node[font=\ttfamily\footnotesize,cGreen] at (0.6,0){15};
    \node[font=\ttfamily\footnotesize,cGreen] at (2.0,0){21};
    \node[font=\ttfamily\footnotesize,cGreen] at (3.4,0){7};
    \node[font=\ttfamily\footnotesize,cGreen] at (4.8,0){3};
    \node[font=\ttfamily\footnotesize,cGreen] at (6.2,0){8};
    % flèches indicatives
    \foreach \x/\i in {0/3,1/2,2/1,3/3,4/0,5/2,6/4}
      \draw[arr,cGray] (v\x) -- (p\i.north);
    \node[font=\scriptsize,cGreen,anchor=west] at (-0.5,-1.5)
      {\textbf{un minimum par partition} $\Rightarrow$ cases indépendantes (comparables case à case)};
  \end{tikzpicture}
\end{frame}

\begin{frame}{Complexité du partition sketch}
  \begin{columns}[T]
    \begin{column}{0.5\textwidth}
      \footnotesize
      \textbf{Construire :} pour chaque k-mer, \texttt{p = kmer \% s} ($O(1)$) puis une
      comparaison $\Rightarrow \mathbf{O(n)}$ au total. Mémoire : $s$ minima $= O(s)$.
      \medskip

      \textbf{Comparer :} parcourir les $s$ slots, \emph{une} comparaison par slot
      $\Rightarrow \mathbf{O(s)}$, \textcolor{cGreen}{sans tri !}
      \[ \hat J = \frac{\#\{i : A_i = B_i\}}{\#\text{slots non vides}} \]
    \end{column}
    \begin{column}{0.5\textwidth}
      \centering
      \begin{tikzpicture}
        \foreach \v/\i in {15/0,21/1,7/2,3/3,8/4} \node[cell] (a\i) at (\i*0.85,0.6){\v};
        \foreach \v/\i in {15/0,21/1,9/2,3/3,8/4} \node[cell] (b\i) at (\i*0.85,-0.6){\v};
        \node[font=\scriptsize,cGray,anchor=east] at (-0.55,0.6){$A$ :};
        \node[font=\scriptsize,cGray,anchor=east] at (-0.55,-0.6){$B$ :};
        \foreach \i in {0,1,3,4} \draw[cGreen,line width=1pt] (a\i.south) -- (b\i.north);
        \draw[cRed,line width=1pt] (a2.south) -- (b2.north);
        \node[cGreen,font=\scriptsize] at (3.6,0.6){\checkmark};
        \node[font=\scriptsize,cGray,align=center] at (1.7,-1.6)
          {4 slots égaux / 5 $\Rightarrow \hat J = 0{,}8$};
      \end{tikzpicture}

      \smallskip
      {\scriptsize
      \begin{tabular}{lcc}
        \toprule
         & construire & comparer \\
        \midrule
        bottom-$s$ & $O(n\log s)$ & $O(s\log s)$ \\
        partition  & $O(n)$       & $O(s)$ \\
        \bottomrule
      \end{tabular}}
    \end{column}
  \end{columns}
\end{frame}

%% ================================================================ PARTIE 4
\section{HyperMin : compression 16 bits}

\begin{frame}{Perte d'espace : les zéros de tête}
  \footnotesize
  On garde des valeurs \emph{parce qu'elles sont petites} : leur écriture binaire
  commence donc par \textbf{beaucoup de zéros}. L'information utile se résume à
  \textbf{(1)} la position du premier \texttt{1} et \textbf{(2)} quelques bits de poids faible.
  \medskip

  \centering
  \begin{tikzpicture}
    \fill[cGray!30] (0,0) rectangle (4.3,0.7);
    \draw (0,0) rectangle (10,0.7);
    \node[font=\ttfamily\footnotesize] at (1.5,0.35){0 0 0 0 0 ... 0};
    \draw[fill=black] (4.3,0) rectangle (4.65,0.7);
    \draw[arr] (4.47,1.35)--(4.47,0.74);
    \node[above,font=\scriptsize] at (4.47,1.3){position du 1\textsuperscript{er} 1 \;($\rho$)};
    \fill[cOrange!45] (8.3,0) rectangle (10,0.7);
    \draw[decorate,decoration={brace,amplitude=4pt},cGray]
      (0,0.8)--(4.3,0.8) node[midway,above,font=\scriptsize,cGray]{zéros de tête = gaspillés};
    \draw[decorate,decoration={brace,amplitude=4pt},cOrange]
      (8.3,0.8)--(10,0.8) node[midway,above,font=\scriptsize,cOrange]{10 bits faibles};
    \node[below,font=\scriptsize] at (0,0){0};
    \node[below,font=\scriptsize] at (10,0){63};
  \end{tikzpicture}

  \medskip
  \footnotesize\textcolor{cGray}{$\Rightarrow$ stocker 64 bits par minimum est un gâchis :
  on peut résumer l'essentiel sur bien moins de bits.}
\end{frame}

\begin{frame}{HyperMin : la signature d'un minimum sur 16 bits}
  \footnotesize
  Un minimum gardé est \textbf{petit} : ses bits de \textbf{poids fort} sont surtout des 0,
  ses bits de \textbf{poids faible} restent aléatoires. On le résume par deux infos.

  \begin{center}
  \begin{tikzpicture}[scale=0.82]
    % --- le minimum (64 bits) ---
    \draw (0,1.4) rectangle (11.4,1.95);
    \fill[cGray!30] (0,1.4) rectangle (4.6,1.95);
    \node[font=\ttfamily\scriptsize] at (1.6,1.675){0 0 0 0 \ldots 0};
    \draw[fill=black] (4.6,1.4) rectangle (4.95,1.95);
    \fill[cOrange!50] (9.5,1.4) rectangle (11.4,1.95);
    \draw[decorate,decoration={brace,amplitude=3pt},cGray] (0,2.0)--(4.5,2.0)
      node[midway,above,font=\scriptsize,cGray]{poids fort (sous pression)};
    \draw[decorate,decoration={brace,amplitude=3pt},cOrange] (9.5,2.0)--(11.4,2.0)
      node[midway,above,font=\scriptsize,cOrange]{poids faible (aléatoires)};
    \draw[arr] (4.77,3.05)--(4.77,1.97);
    \node[above,font=\scriptsize] at (4.77,3.0){le 1 de poids le plus fort};
    % --- la signature (16 bits) : barre proportionnelle 6:10, légendes en dessous ---
    \fill[cBlue!30]   (2.6,-0.55) rectangle (5.0,-0.05);
    \draw[cBlue]      (2.6,-0.55) rectangle (5.0,-0.05);
    \fill[cOrange!45] (5.0,-0.55) rectangle (9.0,-0.05);
    \draw[cOrange]    (5.0,-0.55) rectangle (9.0,-0.05);
    \node[font=\scriptsize,anchor=west] at (9.2,-0.3){$=$ 16 bits};
    \node[font=\scriptsize,cBlue,align=center,below] at (3.8,-0.62){\textbf{6 bits} : position du 1 fort\\($\log_2 64 = 6$)};
    \node[font=\scriptsize,cOrange,align=center,below] at (7.0,-0.62){\textbf{10 bits} :\\poids faible};
    \draw[arr,cBlue]   (4.77,1.4) to[out=-90,in=90] (3.8,-0.03);
    \draw[arr,cOrange] (10.45,1.4) to[out=-90,in=90] (7.0,-0.03);
  \end{tikzpicture}
  \end{center}

  \textcolor{cGray}{Prendre le \textbf{minimum} contraint les bits de poids fort ($\to$ des 0) :
  seule compte \emph{la position du 1 de poids fort}. Les bits de poids faible, plus
  aléatoires, forment la \textbf{signature} discriminante.}
\end{frame}

\begin{frame}{Espace \& compromis mémoire / qualité}
  \begin{columns}[T]
    \begin{column}{0.5\textwidth}
      \footnotesize
      \textbf{Mémoire :}
      \begin{itemize}
        \item Partition : $s\times 64$ bits $= 8s$ octets
        \item HyperMin : $s\times 16$ bits $= 2s$ octets \;\textcolor{cGreen}{($\div 4$)}
      \end{itemize}
      \medskip
      \textbf{Deux régimes :}
      \begin{itemize}
        \item à \textbf{taille égale} : précision quasi identique, 4$\times$ moins de mémoire ;
        \item à \textbf{mémoire égale} : 4$\times$ plus de partitions $\Rightarrow$ écart-type réduit.
      \end{itemize}
      \smallskip
      \textcolor{cRed}{Rançon :} compresser $64\to16$ bits crée des \textbf{collisions}
      $\Rightarrow$ légère \emph{surestimation} de $J$ ; négligeable si les partitions
      sont bien remplies.
    \end{column}
    \begin{column}{0.5\textwidth}
      \centering
      \begin{tikzpicture}
        \node[font=\scriptsize,anchor=east] at (0,1.4){Partition $s$};
        \fill[cBlue!50] (0.1,1.2) rectangle (4.1,1.6);
        \node[font=\tiny,anchor=west] at (4.2,1.4){$8s$ o};
        \node[font=\scriptsize,anchor=east] at (0,0.6){HyperMin $s$};
        \fill[cGreen!55] (0.1,0.4) rectangle (1.1,0.8);
        \node[font=\tiny,anchor=west] at (1.2,0.6){$2s$ o \;($\div4$)};
        \node[font=\scriptsize,anchor=east] at (0,-0.2){HyperMin $4s$};
        \fill[cOrange!55] (0.1,-0.4) rectangle (4.1,0.0);
        \node[font=\tiny,anchor=west] at (4.2,-0.2){$8s$ o};
        \node[font=\scriptsize,cGray,align=center] at (2.2,-1.2)
          {à mémoire égale : 4$\times$ plus de\\partitions $\to$ estimation plus stable};
      \end{tikzpicture}
    \end{column}
  \end{columns}
\end{frame}

%% ================================================================ RECAP
\section{Récapitulatif}

\begin{frame}{Récapitulatif des quatre sketches}
  \centering
  \small
  \begin{tabular}{lcccc}
    \toprule
    \textbf{Sketch} & \textbf{Construire} & \textbf{Comparer} & \textbf{Espace} & \textbf{Tri ?} \\
    \midrule
    Bottom-$s$ (\texttt{set}) & $O(n\,s)$      & $O(s\log s)$ & $64s$ b & oui \\
    Bottom-$s$ (tas)          & $O(n\log s)$   & $O(s\log s)$ & $64s$ b & oui \\
    Partition                 & $O(n)$         & $O(s)$       & $64s$ b & non \\
    \rowcolor{cGreen!15}
    HyperMin                  & $O(n)$         & $O(s)$       & $\mathbf{16s}$ b & non \\
    \bottomrule
  \end{tabular}

  \vfill
  \begin{columns}
    \begin{column}{0.85\textwidth}
      \footnotesize
      \begin{itemize}
        \item \textbf{Encoder en entiers} : énumération et comparaison passent de $O(k)$ à $O(1)$.
        \item \textbf{Bottom-$s$} : échantillon uniforme, mais comparaison liée à un ordre global.
        \item \textbf{Partition} : slots indépendants $\Rightarrow$ comparaison $O(s)$ sans tri.
        \item \textbf{HyperMin} : $\div 4$ de mémoire pour une qualité quasi équivalente.
      \end{itemize}
    \end{column}
  \end{columns}
\end{frame}

\begin{frame}[standout]
  Des k-mers bruts au sketch compact :\\
  \alert{moins de mémoire}, \alert{comparaison plus rapide},\\
  pour une estimation de Jaccard fiable.
\end{frame}

\end{document}
